\documentclass[11pt]{article}

% --- NeurIPS style: use official if present, else stub -------------------
\IfFileExists{neurips_2026.sty}{%
  \usepackage{neurips_2026}%
}{%
  \usepackage{neurips_2026_stub}%
}

% --- standard packages ----------------------------------------------------
\usepackage{amsmath,amssymb,amsthm}
\usepackage{mathtools}
\usepackage{algorithm}
\usepackage{algorithmic}
\usepackage{hyperref}
\usepackage{cleveref}
\usepackage{booktabs}
\usepackage{enumitem}
\usepackage{microtype}
\usepackage{xcolor}

% --- theorem styles -------------------------------------------------------
\newtheorem{proposition}{Proposition}
\newtheorem{remark}{Remark}

% --- title ----------------------------------------------------------------
\title{Fisher-Calibrated Token UCB for Uncertainty-Aware Decoding}
\author{Anonymous Author(s)}

\begin{document}
\maketitle

\begin{abstract}
We formulate uncertainty-aware next-token selection as a local Bayesian
decision problem for a probabilistic language model.
The key object is not the softmax distribution itself, but epistemic
uncertainty in token scores induced by uncertainty in the model parameters.
We approximate the local posterior by a Gaussian around a reference
parameter vector and use the true generalized Gauss--Newton matrix,
which for softmax likelihood coincides with the model Fisher information,
as the posterior precision surrogate.
This yields a token-level upper-confidence decoding rule:
select the token with the largest posterior upper confidence bound on
its log-probability.
We emphasize that this is a decoding rule rather than a contextual-bandit
regret algorithm unless external rewards or preferences are observed online.
\end{abstract}

% --------------------------------------------------------------------------
% 1. Problem and Method
% --------------------------------------------------------------------------
\section{Problem Setup and Method}

\paragraph{Goal.}
At generation step $t$, the model observes a prefix (context) $x_t$
and must choose the next token from a candidate set $\mathcal{A}_t$,
typically the full vocabulary or a computationally truncated shortlist.
A standard decoder chooses according to the point estimate
$p_{\hat\theta}(y \mid x_t)$, for example by argmax, temperature scaling,
top-$k$, or nucleus sampling.
These rules manipulate the model's predictive distribution at fixed
parameters $\hat\theta$, but do not account for epistemic uncertainty
about $\theta$ itself.

We instead treat next-token choice as a local Bayesian decision problem.
For each candidate token $k \in \mathcal{A}_t$, define the scalar score
\[
q_\theta(x_t,k) := \log p_\theta(k \mid x_t).
\]
Our objective is to choose a token with a high posterior upper confidence
bound on $q_\theta(x_t,k)$.

\paragraph{Probabilistic language model.}
Let $z_\theta(x_t) \in \mathbb{R}^{V}$ denote the logits over the
vocabulary of size $V$, and let
\[
p_\theta(k \mid x_t)
=
\frac{\exp(z_\theta(x_t)_k)}
{\sum_{j=1}^{V}\exp(z_\theta(x_t)_j)}.
\]
We write
\[
p_t := p_{\theta_t}(\cdot \mid x_t)
\qquad\text{and}\qquad
q_t(k) := q_{\theta_t}(x_t,k).
\]

\paragraph{Posterior approximation.}
Let $\hat\theta_t$ be a reference parameter vector at generation step $t$.
This may be:
\begin{enumerate}[leftmargin=1.5em]
\item a fixed pretrained checkpoint, with no online parameter updates; or
\item an online-updated parameter vector if feedback is available.
\end{enumerate}
We approximate the local posterior by a Gaussian
\[
p(\theta \mid \mathcal{D}_t)
\approx
\mathcal{N}(\hat\theta_t,\Lambda_t^{-1}),
\]
where $\Lambda_t$ is an approximate posterior precision matrix.

\paragraph{Token-score gradients.}
For each candidate token $k$, define the score gradient
\[
g_t(k) := \nabla_\theta q_{\theta_t}(x_t,k)
=
\nabla_\theta \log p_{\theta_t}(k \mid x_t).
\]
A first-order expansion gives
\[
q_\theta(x_t,k)
\approx
q_t(k) + g_t(k)^\top(\theta-\hat\theta_t).
\]
Under the local Gaussian posterior, the induced posterior variance of
the token score is therefore
\[
s_t^2(k)
:=
g_t(k)^\top \Lambda_t^{-1} g_t(k).
\]

\paragraph{Why the true GGN / Fisher.}
For next-token likelihood, the relevant uncertainty object is the local
curvature of the negative log-likelihood, not a regression Gram matrix
for a scalar reward predictor.
Let
\[
J_t := \frac{\partial z_\theta(x_t)}{\partial \theta}\bigg|_{\theta=\hat\theta_t}
\in \mathbb{R}^{V \times d}.
\]
The true generalized Gauss--Newton matrix for softmax likelihood is
\[
F(x_t;\hat\theta_t)
=
J_t^\top
\Bigl(
\mathrm{Diag}(p_t)-p_t p_t^\top
\Bigr)
J_t.
\]
For softmax cross-entropy, this matrix coincides with the model Fisher
information:
\[
F(x_t;\hat\theta_t)
=
\mathbb{E}_{Y\sim p_{\hat\theta_t}(\cdot\mid x_t)}
\Bigl[
\nabla_\theta \log p_{\hat\theta_t}(Y\mid x_t)\,
\nabla_\theta \log p_{\hat\theta_t}(Y\mid x_t)^\top
\Bigr].
\]
This is the curvature object we use.

\paragraph{Precision matrix.}
We consider two modes.

\emph{Offline decoding mode.}
Fit a local posterior once around a pretrained model using a corpus
$\mathcal{D}_{\mathrm{fit}}$:
\[
\Lambda
=
\lambda I
+
\sum_{(x,y)\in \mathcal{D}_{\mathrm{fit}}}
F(x;\hat\theta).
\]
Then use this fixed precision matrix for all future decoding steps.

\emph{Online update mode.}
If contexts arrive sequentially and the posterior is refreshed online,
define
\[
\Lambda_t
=
\lambda I + \sum_{s<t} F(x_s;\hat\theta_s).
\]
This is the natural sequential analogue.

\paragraph{Token UCB rule.}
The proposed decoding score is
\[
\mathrm{UCB}_t(k)
=
q_t(k) + \beta_t s_t(k),
\]
and the chosen token is
\[
k_t
=
\arg\max_{k\in\mathcal{A}_t}
\mathrm{UCB}_t(k).
\]
The scalar $\beta_t$ sets the confidence level.

\begin{algorithm}[h]
\caption{Fisher-Calibrated Token UCB (FT-UCB)}\label{alg:ftucb}
\begin{algorithmic}[1]
\REQUIRE Language model $p_\theta(\cdot\mid x)$, reference parameters $\hat\theta_t$,
         precision matrix $\hat\Lambda_t$ (fixed or online-refreshed),
         confidence schedule $\{\beta_t\}$,
         candidate-set rule $\mathcal{A}_t$,
         linear solver budget $I$ (or tolerance $\bar\varepsilon$)
\FOR{$t=1,\dots,T$}
  \STATE Observe current prefix $x_t$
  \STATE Form candidate set $\mathcal{A}_t$ (full vocabulary or shortlist)
  \STATE Compute logits $z_{\hat\theta_t}(x_t)$ and probabilities
         $p_t = p_{\hat\theta_t}(\cdot\mid x_t)$
  \FOR{$k\in\mathcal{A}_t$}
    \STATE $g_t(k)\gets \nabla_\theta \log p_{\hat\theta_t}(k\mid x_t)$
    \STATE Solve $\hat\Lambda_t u = g_t(k)$ approximately by CG
    \STATE $\hat s_t^2(k)\gets g_t(k)^\top u$
    \STATE $\mathrm{UCB}_t(k)\gets \log p_{\hat\theta_t}(k\mid x_t)+\beta_t \hat s_t(k)$
  \ENDFOR
  \STATE Select $k_t = \arg\max_{k\in\mathcal{A}_t}\mathrm{UCB}_t(k)$
  \STATE Append token $k_t$ to the prefix
  \STATE Optionally refresh $\hat\theta_{t+1}$ and $\hat\Lambda_{t+1}$ if online updates are used
\ENDFOR
\end{algorithmic}
\end{algorithm}

% --------------------------------------------------------------------------
% 2. Posterior-UCB justification
% --------------------------------------------------------------------------
\section{Posterior-UCB Justification}

The decoding rule above is justified by a local Gaussian posterior and
first-order linearization.

For a fixed context $x_t$ and token $k$, define the Taylor remainder
\[
R_t(k;\theta)
:=
q_\theta(x_t,k)
-
q_t(k)
-
g_t(k)^\top(\theta-\hat\theta_t).
\]
Let the local ellipsoid be
\[
\mathcal{E}_t(\beta_t)
:=
\left\{
\theta:
(\theta-\hat\theta_t)^\top
\Lambda_t
(\theta-\hat\theta_t)
\le \beta_t^2
\right\}.
\]

\begin{proposition}[Token-wise posterior confidence bound]
Assume
\[
\theta \mid \mathcal{D}_t \sim \mathcal{N}(\hat\theta_t,\Lambda_t^{-1}),
\]
and suppose that for every $k\in\mathcal{A}_t$ and every
$\theta\in\mathcal{E}_t(\beta_t)$,
\[
|R_t(k;\theta)| \le \varepsilon_{\mathrm{lin}}(t,k).
\]
Then for any $\delta_t\in(0,1)$, if
\[
\beta_t = \sqrt{2\log\!\left(\frac{2|\mathcal{A}_t|}{\delta_t}\right)},
\]
we have, with posterior probability at least $1-\delta_t$,
simultaneously for all $k\in\mathcal{A}_t$,
\[
q_\theta(x_t,k)
\le
q_t(k)+\beta_t s_t(k)+\varepsilon_{\mathrm{lin}}(t,k),
\]
and
\[
q_\theta(x_t,k)
\ge
q_t(k)-\beta_t s_t(k)-\varepsilon_{\mathrm{lin}}(t,k).
\]
\end{proposition}

\begin{proof}
Under the Gaussian posterior,
\[
g_t(k)^\top(\theta-\hat\theta_t)
\sim
\mathcal{N}(0,s_t^2(k)).
\]
Hence, for each fixed $k$,
\[
\mathbb{P}\!\left(
\left|g_t(k)^\top(\theta-\hat\theta_t)\right|
>
\beta_t s_t(k)
\right)
\le
2e^{-\beta_t^2/2}.
\]
Applying a union bound over $k\in\mathcal{A}_t$ gives simultaneous control
with probability at least
\[
1 - 2|\mathcal{A}_t|e^{-\beta_t^2/2}.
\]
Choosing
\[
\beta_t = \sqrt{2\log\!\left(\frac{2|\mathcal{A}_t|}{\delta_t}\right)}
\]
makes this failure probability at most $\delta_t$.
Adding the deterministic linearization remainder bound yields the claim.
\end{proof}

\paragraph{Interpretation.}
The score
\[
q_t(k)+\beta_t s_t(k)
\]
is therefore a posterior upper confidence bound on the token
log-probability, up to the local linearization error
$\varepsilon_{\mathrm{lin}}(t,k)$.
This is the correct analogue of UCB for token selection under local
parameter uncertainty.

% --------------------------------------------------------------------------
% 3. Which curvature matrix should be used?
% --------------------------------------------------------------------------
\section{Which Curvature Matrix Should Be Used?}

\paragraph{Observed Hessian.}
For a realized token $y$, the negative log-likelihood Hessian is
\[
H_{\mathrm{obs}}(x,y;\theta)
=
\nabla_\theta^2\bigl[-\log p_\theta(y\mid x)\bigr].
\]
This is the exact local curvature of the observed loss.
In principle, a Laplace approximation could use
\[
\Lambda \approx \lambda I + \sum H_{\mathrm{obs}}(x,y;\hat\theta).
\]
However, this matrix contains second-derivative terms of the network
mapping itself and is often less stable computationally.

\paragraph{True GGN / model Fisher.}
For softmax likelihood, the true GGN is
\[
F(x;\theta)
=
J^\top\bigl(\mathrm{Diag}(p)-pp^\top\bigr)J,
\]
which is positive semidefinite and coincides with the model Fisher.
It is label-averaged, tied directly to the predictive distribution,
and is therefore the preferred curvature matrix for token-level
epistemic uncertainty.

\paragraph{Empirical Fisher.}
The empirical Fisher for a realized token $y$ is
\[
F_{\mathrm{emp}}(x,y;\theta)
=
\nabla_\theta \log p_\theta(y\mid x)\,
\nabla_\theta \log p_\theta(y\mid x)^\top.
\]
We do \emph{not} use this object for uncertainty-aware decoding.
It depends on the realized label $y$, whereas at test time the next token
has not yet been observed.
Moreover, it is not the same as the model Fisher in general.

\begin{remark}[Recommended choice]
For next-token decoding under a softmax language model, the correct
default curvature matrix is the true GGN / model Fisher.
If one is fitting a scalar reward model instead, then the appropriate
uncertainty matrix is the Jacobian Gram of that reward predictor, not the
token-likelihood Fisher.
\end{remark}

% --------------------------------------------------------------------------
% 4. Practical approximations
% --------------------------------------------------------------------------
\section{Practical Approximations}

\paragraph{Shortlists.}
For large vocabularies, computing $g_t(k)$ and $s_t(k)$ for every token
is expensive.
In practice one may restrict to a shortlist
\[
\mathcal{A}_t \subseteq \{1,\dots,V\},
\]
for example the top-$M$ tokens by $\log p_{\hat\theta_t}(k\mid x_t)$.

\paragraph{Matrix-free solves.}
The variance term requires solves of the form
\[
\Lambda_t u = g_t(k).
\]
These can be computed by conjugate gradients using only matrix-vector
products with $\Lambda_t$.
When $\Lambda_t$ is built from the true GGN / Fisher, these products can
be implemented via Jacobian-vector and vector-Jacobian operations.

\paragraph{Fixed-posterior decoding.}
If the model is not updated online, one can fit a single local posterior
around a reference checkpoint and use the same $\Lambda$ throughout
decoding.
In that case the method is a test-time uncertainty-aware decoder.

\paragraph{Online updates.}
If external feedback is available, such as token rewards, preferences,
or sequence-level outcomes, one may update $\hat\theta_t$ and $\Lambda_t$
online.
That setting is closer to a genuine bandit or reinforcement-learning
problem, but is distinct from pure decoding.

% --------------------------------------------------------------------------
% 5. Experimental protocol
% --------------------------------------------------------------------------
\section{Experimental Protocol}

We propose evaluating Fisher-Calibrated Token UCB in two regimes.

\paragraph{Offline decoding.}
Take a pretrained language model and fit a local Gaussian posterior
around the checkpoint using a held-out fitting corpus.
Compare FT-UCB decoding against greedy decoding, temperature sampling,
top-$k$, and nucleus sampling.
Metrics include negative log-likelihood, exact-match accuracy on next-token
prediction, calibration of posterior intervals for token scores,
sequence-level task performance, and diversity.

\paragraph{Online feedback.}
In settings with external rewards or preferences, compare FT-UCB against
standard exploration baselines.
In this regime, one may update $\hat\theta_t$ and $\Lambda_t$ online and
measure cumulative reward or regret.

\paragraph{Ablations.}
Key ablations include:
\begin{enumerate}[leftmargin=1.5em]
\item full vocabulary versus shortlist decoding;
\item last-layer versus full-model curvature;
\item exact versus approximate linear solves;
\item fixed posterior versus online posterior refresh;
\item true GGN / Fisher versus observed Hessian approximations.
\end{enumerate}

% --------------------------------------------------------------------------
% 6. Scope and limitation
% --------------------------------------------------------------------------
\section{Scope and Limitation}

This method provides a principled way to use local parameter uncertainty
for token selection.
However, in the absence of external feedback, it should be interpreted as
an uncertainty-aware decoding rule, not as a contextual-bandit algorithm
with a regret guarantee.
A regret theorem would require a separate reward model and an explicit
feedback process.

% --------------------------------------------------------------------------
% 7. Conclusion
% --------------------------------------------------------------------------
\section{Conclusion}

For next-token selection under a probabilistic language model, the correct
uncertainty object is the posterior uncertainty of token scores induced by
parameter uncertainty.
Under a local Gaussian approximation, this leads to a token-level UCB rule
whose covariance is determined by the inverse posterior precision.
For softmax likelihood, the natural precision surrogate is the true
generalized Gauss--Newton matrix, equivalently the model Fisher
information.
This yields a clean and principled alternative to heuristic decoding rules.

\end{document}